% -------------------------------------------------------
%  Section 1: Paper Identity & Executive Summary
% -------------------------------------------------------
\section{شناسنامه مقاله و چکیده مدیریتی}
\label{sec:identity}

مقاله مورد بررسی در جدول \ref{tab:bibinfo} معرفی شده است. این اثر یک پیش‌چاپ \lr{arXiv} است که برای ارسال به \lr{Elsevier} آماده شده و داوری همتا را پشت سر نگذاشته؛ همین نکته باید در وزن‌دهی به ادعا‌های آن لحاظ شود. پانویس صفحه نخست نشانی یک مخزن عمومی را برای کد و داده اعلام می‌کند، ولی آن نشانی در دسترس نیست و مخزنی با آن نام در حساب نویسنده وجود ندارد. بنابراین مصنوع پژوهشی این مقاله عملاً منتشرنشده است و هر بازتولیدی ناگزیر باید از روی متن انجام شود؛ همان کاری که در بخش دوم این گزارش کرده‌ایم.

\begin{table}[!ht]
    \centering
    \footnotesize
    \caption{شناسنامه کتاب‌شناختی مقاله مورد بررسی}
    \label{tab:bibinfo}
    \begin{tabular}{|p{2.6cm}|p{11.4cm}|}
    \hline
    \textbf{عنوان} & \lr{Malware Classification Based on Image Segmentation} \\ \hline
    \textbf{نویسنده} & \lr{Wanhu Nie} (نویسنده مسئول و تنها نویسنده)، دانشگاه صنعتی لانژو، چین \\ \hline
    \textbf{نوع و سال} & پیش‌چاپ \lr{arXiv:2406.03831v1}، ۶ ژوئن ۲۰۲۴، هشت صفحه، ارسال‌شده به \lr{Elsevier} \\ \hline
    \textbf{کلیدواژه‌ها} & دسته‌بندی بدافزار، مصورسازی بدافزار\pn{Malware Visualization}، تحلیل ایستا\pn{Static Analysis}، یادگیری عمیق \\ \hline
    \textbf{مجموعه‌داده} & \lr{Microsoft Malware Classification Challenge (BIG 2015)} \cite{ronen2018big2015} \\ \hline
    \textbf{دسترسی به کد} & نشانی مخزن در پانویس مقاله آمده است ولی در دسترس نیست (بررسی‌شده در شهریور ۱۴۰۵) \\ \hline
    \end{tabular}
\end{table}

روش‌های «مصورسازی بدافزار» جریان بایتی یک فایل اجرایی را به‌عنوان شدت روشنایی پیکسل تفسیر می‌کنند و مسئله دسته‌بندی بدافزار را به دسته‌بندی تصویر فرو می‌کاهند \cite{nataraj2011images}. نقطه ضعف شناخته‌شده این خانواده آن است که مهاجم می‌تواند با جابه‌جایی بخش‌ها\pn{Sections} یا تزریق داده زائد، چیدمان فضایی تصویر را دستکاری کند و طبقه‌بند را گمراه سازد.

ایده مرکزی مقاله ساده و روشن است: به‌جای نگاشت کل فایل به یک تصویر واحد، تصویر خاکستری بر اساس \textbf{نوع بخش} قطعه‌بندی شود و هر زیرتصویر به‌عنوان یک \textbf{کانال} جداگانه به شبکه پیچشی برود. دو صورت‌بندی پیشنهاد می‌شود: \lr{Split} که هر بخش را فشرده و مستقل از جایگاهش ذخیره می‌کند، و \lr{Mask} که موقعیت اصلی بخش را حفظ ولی بقیه تصویر را صفر می‌کند. در کنار این، سه طرح هم‌ترازی عرض مقایسه و عرض تصویر به‌عنوان یک ابرپارامتر مستقل معرفی می‌شود. ارزیابی روی \lr{BIG 2015} با \lr{VGG16} \cite{simonyan2015vgg} و \lr{ResNet50} \cite{he2016resnet} و ۲۰ پیکربندی انجام شده و سنجه اصلی، زیان لگاریتمی چندکلاسه\pn{Multi-class Logarithmic Loss} سرور ارزیابی کگل است؛ بهترین عدد گزارش‌شده \num{۰٫۰۲۶۵} برای \lr{ResNet50} روی \lr{S3} (عرض ثابت \num{۱۰۲۴}) است.

این گزارش دو بخش دارد. بخش نخست (\S\ref{sec:identity} تا \S\ref{sec:critique}) مقاله را تحلیل و نقد می‌کند و هر ایراد را تنها بر پایه محتوای منتشرشده خود مقاله مستند می‌سازد. بخش دوم (\S\ref{sec:repro} تا \S\ref{sec:conclusion}) کار تیم ما در فازهای سوم و چهارم را گزارش می‌کند: بازتولید کامل ۲۰ پیکربندی مقاله روی \lr{BIG 2015}، انتقال روش به دامنه سیستم‌های نهفته با یک مجموعه‌داده \lr{ELF} از \lr{Zephyr} روی \lr{ARM}، و یک مطالعه حذفی\pn{Ablation Study} روی مجموعه‌ای از بهبود‌ها. در بخش دوم همان سخت‌گیری را به کار خودمان نیز اعمال کرده‌ایم و نشان می‌دهیم پروتکل ارزیابی اولیه‌مان معیوب بوده و اعداد اصلاح‌شده چه تصویر متفاوتی می‌سازند.
